% =============================================================================
% Textbausteine fuer haeufige Abkuerzungen
% =============================================================================
% Warum das wichtig ist: "z. B." wird mit einem SCHMALEN, GESCHUETZTEN
% Leerzeichen gesetzt. Schreibst du "z. B." mit normalem Leerzeichen, darf
% LaTeX zwischen "z." und "B." umbrechen -- ein klassischer Typografiefehler,
% der in der Korrektur auffaellt. Schreibst du "z.B." ohne Leerzeichen, ist es
% ebenfalls falsch.
%
% Benutzung im Fliesstext:   \zb  \dah  \ua  \vgl  \bzw  \ca
% Vor einem Komma oder Punkt die OL-Variante (ohne folgendes Leerzeichen):
%   ... \zbol, ...
%
% Uebernommen aus dem FOM-LaTeX-Template von Andy Grunwald (MIT) und auf die
% an der FOM gebraeuchlichen Formen reduziert.
% =============================================================================

\newcommand{\zbol}{z.\,B.}
\newcommand{\zb}{\zbol\ }
\newcommand{\dahol}{d.\,h.}  % \dh ist von LaTeX belegt (Eth-Zeichen), daher \dah
\newcommand{\dah}{\dahol\ }
\newcommand{\uaol}{u.\,a.}
\newcommand{\ua}{\uaol\ }
\newcommand{\vglol}{Vgl.}
\newcommand{\vgl}{\vglol\ }
\newcommand{\bzwol}{bzw.}
\newcommand{\bzw}{\bzwol\ }
\newcommand{\caol}{ca.}
\newcommand{\ca}{\caol\ }
\newcommand{\ggfol}{ggf.}
\newcommand{\ggf}{\ggfol\ }
\newcommand{\os}{o.\,S.}      % ohne Seitenangabe (Online-Quellen)
\newcommand{\oj}{o.\,J.}      % ohne Jahr
\newcommand{\ov}{o.\,V.}      % ohne Verfasser
